\section{Тест производительности}

Для оценки эффективности реализованного АВЛ-дерева было проведено сравнение со стандартной хэш-таблицей языка C++ --- \texttt{std::unordered\_map} (стандарт C++20). 

Обе структуры данных тестировались на трех наборах данных (по 1,000,000 операций в каждом):
\begin{enumerate}
    \item \textbf{Random Insertion:} вставка случайно сгенерированных строк.
    \item \textbf{Sorted Insertion:} вставка лексикографически отсортированных строк.
    \item \textbf{Search:} поиск существующих и несуществующих элементов в заполненном словаре.
\end{enumerate}

Измерения проводились с помощью утилиты \texttt{time} (метрика \texttt{user time}, отражающая чистое процессорное время). Результаты представлены в таблице \ref{tab:bench_results}.

\begin{table}[h!]
\centering
\begin{tabular}{|l|c|c|c|}
\hline
\textbf{Сценарий (1 млн оп.)} & \textbf{AVL Tree (user)} & \textbf{std::unordered\_map} & \textbf{Ускорение} \\ \hline
Random Insertion & 1.207 с & \textbf{0.453 с} & Map в 2.6x \\ \hline
Sorted Insertion & \textbf{0.295 с} & 0.440 с & AVL в 1.5x \\ \hline
Search           & 0.336 с & \textbf{0.113 с} & Map в 2.9x \\ \hline
\end{tabular}
\caption{Сравнение времени работы АВЛ-дерева и хэш-таблицы}
\label{tab:bench_results}
\end{table}

\subsection*{Анализ результатов}

1. \textbf{Случайные данные и поиск.} На случайных данных и при поиске хэш-таблица ожидаемо побеждает. Асимптотика \texttt{std::unordered\_map} в среднем составляет $O(1)$, в то время как АВЛ-дерево работает за $O(\log N)$. Более того, рекурсивная реализация дерева с использованием \texttt{std::unique\_ptr} накладывает значительные расходы на вызовы методов умных указателей при спуске по дереву.

2. \textbf{Отсортированные данные.} На отсортированных данных АВЛ-дерево оказалось \textbf{в 1.5 раза быстрее} хэш-таблицы. Это связано с тем, что при вставке лексикографически возрастающих ключей хэш-таблица вынуждена постоянно перевыделять память (rehashing), пересчитывая хэши для всех ранее вставленных элементов. АВЛ-дерево в этом сценарии выполняет предсказуемые балансировки (малые левые вращения) и всегда идет в правую ветку, что отлично ложится в кэш процессора.

\pagebreak